%
%      Brno University of Technology
%      Faculty of Information Technology
%
%      MSc Thesis	2009/2010
%
%      Design of S-Boxes Using Genetic Algorithms
%
%
%      Bedrich Hovorka
%
%
Prvním výsledkem diplomové práce je implementovaný modul s knihovními funkcemi pro Python, který
využívá algoritmy implementované v dynamicky linkované knihovně napsané v jazyce C++.
Druhým výsledkem jsou experimentálně zjištěné hodnoty.
Bylo navrženo jednoduché rozhraní pro popis parametrů evoluce, aby ho bylo možné
v budoucnu použít pro další experimenty.
Ve fázi implementace bylo nutné efektivně naprogramovat různé evoluční algoritmy.
K tomu byly využity nástroje Valgrind, který pomáhá odhalit chyby při práci s pamětí, a OProfile, jenž hledá části kódu, které
trvají nejdéle.
Aby Python rozhraní umožňovalo jednoduše popsat parametry
evoluce, bylo nutné provést úpravy na úrovni C++, což se projevilo v jednodušším popisu
experimentů a zpracování výsledků.
Experimenty ukázaly, které algoritmy a jejich parametry jsou vhodnější pro určitá
kritéria na zvolených reprezentacích, a to jak
pro jednotlivá kritéria bezpečnosti substitučních boxů, tak i pro multikriteriální
optimalizaci.
V práci \cite{evo} je zmíněno, že bylo provedeno jen málo srovnání
genetických algoritmů a algoritmů EDA.
I z tohoto důvodu bylo provedeno srovnání na tomto konkrétním problému.
Nebylo možné otestovat všechny kombinace algoritmů, reprezentací, kritérií a parametrů.
K experimentům byly vybrány ty, které se zdály být nejzajímavější.

\subsection*{Náměty pro další rozvoj}
Pro S-boxy složitější, než se kterými se experimentovalo v této práci, je již
nutné zvážit hardwarovou akceleraci výpočtů jednotlivých kritérií a prohledávacích
algoritmů.
Dalším kritériem by mohl být řád polynomu v $GF(2^m)$.
Vyšší řád polynomu \cite{langrangeAttack, courtois} znamená vyšší odolnost vůči útoku pomocí interpolace Lagrangeovým
polynomem.
Jedná se o zobecnění lineární kryptoanalýzy. Právě pomocí Lagrangeovy interpolace
se získá polynom vyjadřující S-box. Jeho stupeň lze dále maximalizovat.
Existuje ovšem několik variant Galoisových polí, ze kterých si útočník může vybrat, tudíž je nutná
odolnost vůči všem.
